\documentclass{amsart}
\usepackage{amsmath,amssymb,amsthm,hyperref}
\newtheorem{theorem}{Theorem}
\newtheorem{lemma}{Lemma}
\newtheorem{proposition}{Proposition}
\newtheorem{corollary}{Corollary}
\theoremstyle{definition}
\newtheorem{definition}{Definition}
\newtheorem{remark}{Remark}
\begin{document}

\title{Connected signed graphs with two distinct eigenvalues}
\maketitle

\section{Setup, conventions and statement of the result}

Throughout, $(G,\sigma)$ is a signed graph with $G=(V,E)$ a finite simple
graph on $n=|V|\ge 2$ vertices and $\sigma\colon E\to\{+1,-1\}$ a signature.
We write $A=A^\sigma$ for the symmetric $n\times n$ adjacency matrix:
$A_{uv}=\sigma(uv)$ if $uv\in E$ and $A_{uv}=0$ otherwise; in particular
$A_{uu}=0$. As $A$ is real symmetric it has $n$ real eigenvalues counted with
multiplicity. We assume $G$ is connected.

\paragraph{Switching.}
For $D=\operatorname{diag}(\varepsilon_1,\dots,\varepsilon_n)$ with
$\varepsilon_i\in\{+1,-1\}$, the matrix $D A D$ is again the adjacency matrix
of a signed graph on the same underlying graph $G$, with signature
$\sigma'(uv)=\varepsilon_u\varepsilon_v\sigma(uv)$. Two signed graphs related
this way are \emph{switching equivalent}. Since $D=D^{-1}=D^\top$, $DAD$ is
similar to $A$, so switching equivalent signed graphs have the same spectrum;
switching also preserves connectedness. We classify up to switching
equivalence and isomorphism, as required.

\paragraph{Notation.} $I$ is the identity matrix, $J$ the all-ones matrix,
$\mathbf 1$ the all-ones vector. We repeatedly use
\begin{equation}\label{eq:walk}
(A^2)_{ij}=\sum_{h=1}^n A_{ih}A_{hj}\quad(i\ne j),\qquad
(A^2)_{ii}=\sum_{h=1}^n A_{ih}^2=\deg_G(i),
\end{equation}
the last equality because $A_{ih}^2=1$ exactly when $ih\in E$ and $0$ otherwise.

\begin{definition}[Symmetric weighing matrix]\label{def:weigh}
A \emph{symmetric weighing matrix of weight $k$ and order $n$ with zero
diagonal} is a symmetric $n\times n$ matrix $W$ with entries in
$\{-1,0,+1\}$, zero diagonal, and $W^2=kI$. Equivalently it is the adjacency
matrix of a signed graph in which every vertex has degree $k$ and $A^2=kI$. We
write $W(n,k)$ for the class of such matrices.
\end{definition}

\begin{theorem}\label{thm:main}
Let $(G,\sigma)$ be a connected signed graph on $n\ge 2$ vertices with
adjacency matrix $A$. Then $A$ has exactly two distinct eigenvalues $a>b$ if
and only if there is an integer $k\ge1$ such that $G$ is $k$-regular and
\begin{equation}\label{eq:quad}
A^2=(a+b)\,A+k\,I,\qquad k=-ab .
\end{equation}
Writing $s=a+b$, exactly one of the following holds.

\smallskip
\textbf{(I) Balanced case $s=0$.}
Then $a=\sqrt k$, $b=-\sqrt k$, both eigenvalues have multiplicity $n/2$ (so
$n$ is even), and \eqref{eq:quad} reads $A^2=kI$. Thus $A\in W(n,k)$ is a
connected symmetric weighing matrix of weight $k$ with zero diagonal;
conversely every connected $A\in W(n,k)$ is the adjacency matrix of a signed
graph with the two eigenvalues $\pm\sqrt k$.

\smallskip
\textbf{(II) Integral case $s\ne 0$.}
Then $a$ and $b$ are integers with $a>0>b$, $a+b=s$, $ab=-k$, and
\eqref{eq:quad} holds with integer coefficients. The two extreme multiplicity
patterns are completely determined:
\begin{itemize}
\item[(IIa)] If $a$ has multiplicity $1$ then, up to switching, $A=J-I$;
$(G,\sigma)$ is switching equivalent to the all-positive $K_n$, with spectrum
$\{\,n-1,\ \underbrace{-1,\dots,-1}_{n-1}\,\}$, i.e.\ $a=n-1$, $b=-1$.
\item[(IIb)] If $b$ has multiplicity $1$ then, up to switching, $A=-(J-I)$;
$(G,\sigma)$ is switching equivalent to the all-negative $K_n$, with spectrum
$\{\,\underbrace{1,\dots,1}_{n-1},\ -(n-1)\,\}$, i.e.\ $a=1$, $b=-(n-1)$.
\end{itemize}
If both multiplicities are at least $2$, then $A$ is a connected, $k$-regular,
zero-diagonal $\{0,\pm1\}$ symmetric matrix satisfying \eqref{eq:quad} with the
integer $s\ne0$; these are exactly the connected integral two-eigenvalue signed
graphs.
\end{theorem}

Section~\ref{sec:quad} proves the equivalence with \eqref{eq:quad} and
regularity; Section~\ref{sec:dichotomy} establishes the dichotomy and the
integrality in case~(II); Section~\ref{sec:simple} settles the simple
eigenvalue subcases; Section~\ref{sec:examples} records the converses and the
nonemptiness of each family. The theorem is a complete characterization: it
states, up to switching and isomorphism, exactly which signed graphs occur,
identifying each family with a precisely defined intrinsic class.

\section{The defining quadratic equation and regularity}\label{sec:quad}

\begin{lemma}\label{lem:quad}
For two distinct reals $a>b$, the matrix $A$ has eigenvalues lying in
$\{a,b\}$ and uses both values if and only if $(A-aI)(A-bI)=0$, equivalently
\begin{equation}\label{eq:quad2}
A^2=(a+b)A-ab\,I .
\end{equation}
In particular $A$ has exactly two distinct eigenvalues if and only if
\eqref{eq:quad2} holds for some reals $a>b$.
\end{lemma}

\begin{proof}
By the Spectral Theorem write $A=\sum_{\lambda}\lambda P_\lambda$, the sum over
the distinct eigenvalues $\lambda$, where $P_\lambda$ is the orthogonal
projection onto the $\lambda$-eigenspace, $\sum_\lambda P_\lambda=I$ and
$P_\lambda P_\mu=0$ for $\lambda\ne\mu$.

If every eigenvalue lies in $\{a,b\}$ then
$(A-aI)(A-bI)=\sum_\lambda(\lambda-a)(\lambda-b)P_\lambda=0$ because each scalar
$(\lambda-a)(\lambda-b)$ vanishes. Conversely, if $(A-aI)(A-bI)=0$ then
applying it to any eigenvector of eigenvalue $\lambda$ gives
$(\lambda-a)(\lambda-b)=0$, so $\lambda\in\{a,b\}$.

Now suppose \eqref{eq:quad2}, i.e.\ $(A-aI)(A-bI)=0$, holds with $a\ne b$. The
distinct eigenvalues form a nonempty subset of $\{a,b\}$; it cannot be a single
value $c$, for then $A=cI$, and the zero diagonal forces $c=0$, hence $A=0$,
contradicting that the connected graph $G$ on $n\ge2$ vertices has an edge
(so $A\ne0$). Hence both $a,b$ occur and $A$ has exactly two distinct
eigenvalues. The displayed identity
$A^2-(a+b)A+abI=(A-aI)(A-bI)$ gives the equivalence of the two forms.
\end{proof}

\begin{lemma}\label{lem:regular}
If $A$ has exactly two distinct eigenvalues $a>b$, then $-ab$ is a positive
integer $k$, $G$ is $k$-regular, $a>0>b$, and \eqref{eq:quad} holds.
\end{lemma}

\begin{proof}
By Lemma~\ref{lem:quad}, \eqref{eq:quad2} holds. Comparing $i$-th diagonal
entries and using \eqref{eq:walk} with $A_{ii}=0$,
\[
\deg_G(i)=(A^2)_{ii}=(a+b)A_{ii}-ab=-ab\qquad\text{for every }i .
\]
Thus all degrees equal $k:=-ab$. As $G$ is connected on $n\ge2$ vertices it has
an edge, so some degree is $\ge1$; hence $k\ge1$ and $k\in\mathbb Z$, since
every degree is a nonnegative integer. From $-ab=k>0$ and $a>b$ we get
$a>0>b$. Substituting $-ab=k$ in \eqref{eq:quad2} yields \eqref{eq:quad}; $G$
is $k$-regular.
\end{proof}

The ``only if'' direction of the equivalence in Theorem~\ref{thm:main} (with
$k$-regularity and $k=-ab\ge1$) is Lemma~\ref{lem:regular}. The ``if''
direction follows from Lemma~\ref{lem:quad}: if $G$ is $k$-regular and
\eqref{eq:quad} holds with $k=-ab$, then $(A-aI)(A-bI)=0$, and $a\ne b$ because
$k=-ab>0$ forces opposite signs; hence $A$ has exactly two distinct
eigenvalues.

\section{The dichotomy $s=0$ versus $s\ne0$, and integrality}\label{sec:dichotomy}

Let $m_a,m_b\ge1$ denote the multiplicities of $a,b$, so $m_a+m_b=n$. Because
$A$ has zero diagonal, $\operatorname{tr}A=0$, hence
\begin{equation}\label{eq:trace}
m_a\,a+m_b\,b=0 .
\end{equation}

\begin{lemma}\label{lem:dichotomy}
Set $s=a+b$. Then either
\begin{enumerate}
\item[(I)] $s=0$, in which case $a=\sqrt k$, $b=-\sqrt k$, $m_a=m_b=n/2$ (so
$n$ is even), and $A^2=kI$; or
\item[(II)] $s\ne0$, in which case $a$ and $b$ are integers.
\end{enumerate}
\end{lemma}

\begin{proof}
\emph{Case $s=0$.} Then $b=-a$ and $a^2=-ab=k$, so $a=\sqrt k$, $b=-\sqrt k$
(recall $a>0$). Equation~\eqref{eq:trace} becomes $a(m_a-m_b)=0$; since $a>0$,
$m_a=m_b$, so $n=2m_a$ is even and each multiplicity is $n/2$. With $a+b=0$,
\eqref{eq:quad} reads $A^2=kI$.

\emph{Case $s\ne0$.} Since all entries of $A$ are integers, its characteristic
polynomial
\[
p(x)=\det(xI-A)=(x-a)^{m_a}(x-b)^{m_b}
\]
has integer coefficients (it is the determinant of a matrix with entries in
$\mathbb Z[x]$). In particular $a$ is a root of a monic integer polynomial, so
$a$ is an algebraic integer.

Suppose, for contradiction, that $a$ is irrational. Let $q\in\mathbb Z[x]$ be
the minimal polynomial of $a$ over $\mathbb Q$; it is monic with integer
coefficients (the minimal polynomial of an algebraic integer is monic in
$\mathbb Z[x]$) and of degree $d\ge2$ since $a\notin\mathbb Q$. As $q$ is the
minimal polynomial of a root of $p$, we have $q\mid p$ in $\mathbb Q[x]$, so
every root of $q$ is a root of $p$ and hence lies in $\{a,b\}$. Over a field of
characteristic $0$, $q$ is separable, so its $d\ge2$ roots are distinct; being
a subset of the two-element set $\{a,b\}$, they are exactly $\{a,b\}$ and
$d=2$. Therefore
\[
q(x)=(x-a)(x-b)=x^2-s\,x-k\in\mathbb Z[x],
\]
so $s\in\mathbb Z$ and $k\in\mathbb Z$, and $a,b$ are the two (irrational)
Galois conjugates of one another.

We now show $m_a=m_b$. Let $\tau\in\operatorname{Gal}(\overline{\mathbb
Q}/\mathbb Q)$ be a field automorphism with $\tau(a)=b$ (such $\tau$ exists
because $a,b$ are conjugate roots of the irreducible $q$). Applying $\tau$ to
the coefficients of $p$ fixes $p$ (its coefficients are rational), so
$\tau$ permutes the roots of $p$ preserving multiplicities; since
$\tau(a)=b$, the multiplicity of $a$ equals the multiplicity of $b$, i.e.\
$m_a=m_b$. Then \eqref{eq:trace} gives $m_a(a+b)=0$, hence $a+b=s=0$,
contradicting $s\ne0$.

Therefore $a$ is rational. Being an algebraic integer and rational, $a$ is an
ordinary integer (a rational algebraic integer lies in $\mathbb Z$). The same
reasoning applies to $b$: it is an eigenvalue of the integer matrix $A$, hence an
algebraic integer, and $b=-k/a\in\mathbb Q$ (with $a\ne0$ since $a>0$), so
$b\in\mathbb Z$. Consequently $a,b\in\mathbb Z$, $s=a+b\in\mathbb Z$ and
$k=-ab\in\mathbb Z$.
\end{proof}

\section{The simple-eigenvalue subcases}\label{sec:simple}

\begin{proposition}\label{prop:simple}
Suppose $A$ has exactly two distinct eigenvalues $a>b$ and $a$ has
multiplicity $1$. Then $b=-1$, $a=n-1$, and $DAD=J-I$ for some $\pm1$ diagonal
$D$; thus $G=K_n$ and $(G,\sigma)$ is switching equivalent to the all-positive
complete graph. Symmetrically, if $b$ has multiplicity $1$, then $a=1$,
$b=-(n-1)$, and $DAD=-(J-I)$ for some $\pm1$ diagonal $D$.
\end{proposition}

\begin{proof}
Assume $m_a=1$ and let $v$ be a unit eigenvector for $a$. By the Spectral
Theorem the projections onto the $a$- and $b$-eigenspaces are $vv^\top$ and
$I-vv^\top$, so
\[
A=a\,vv^\top+b\,(I-vv^\top)=b\,I+(a-b)\,vv^\top .
\]
Comparing diagonal entries with $A_{ii}=0$:
\[
0=b+(a-b)v_i^2\ \Longrightarrow\ v_i^2=\frac{-b}{a-b}=:t,
\]
the same positive value $t$ for all $i$ (positive because $a>0>b$). Since $v$
is a unit vector, $nt=\sum_i v_i^2=1$, so $t=1/n$ and $|v_i|=\sqrt t>0$ for all
$i$.

For $i\ne j$, $A_{ij}=(a-b)v_iv_j$, so
$|A_{ij}|=(a-b)|v_i||v_j|=(a-b)t=(a-b)\cdot\frac{-b}{a-b}=-b>0$. But
$A_{ij}\in\{-1,0,1\}$ is nonzero, so $|A_{ij}|=1$, giving $-b=1$, i.e.\ $b=-1$.
Hence $A_{ij}\ne0$ for all $i\ne j$: the underlying graph is $K_n$, and using
$(a-b)t=-b=1$ and $t=|v_i||v_j|$,
\[
A_{ij}=(a-b)v_iv_j=\frac{v_iv_j}{|v_i||v_j|}
=\operatorname{sign}(v_i)\operatorname{sign}(v_j)\qquad(i\ne j).
\]
Let $D=\operatorname{diag}(\operatorname{sign}(v_1),\dots,
\operatorname{sign}(v_n))$, a $\pm1$ matrix (no $v_i=0$). Then for $i\ne j$,
$(DAD)_{ij}=\operatorname{sign}(v_i)A_{ij}\operatorname{sign}(v_j)=1$ and
$(DAD)_{ii}=0$, so $DAD=J-I$. The all-positive $K_n$ has eigenvalues $n-1$
(simple, eigenvector $\mathbf 1$) and $-1$ (multiplicity $n-1$), so $a=n-1$,
$b=-1$.

If instead $b$ has multiplicity $1$, apply the above to $-A$, the adjacency
matrix of $(G,-\sigma)$, whose largest eigenvalue $-b$ is then simple. We get
$D(-A)D=J-I$, i.e.\ $DAD=-(J-I)$, and the eigenvalues of $-A$ are $n-1$
(simple) and $-1$, so those of $A$ are $a=1$ and $b=-(n-1)$ (simple).
\end{proof}

\section{Converses and nonemptiness of the families}\label{sec:examples}

\paragraph{Case (I), $s=0$.}
If $A\in W(n,k)$ is connected, then $A$ is symmetric with entries in
$\{0,\pm1\}$, zero diagonal, and $A^2=kI$, i.e.\ $(A-\sqrt kI)(A+\sqrt kI)=
A^2-kI=0$; by Lemma~\ref{lem:quad} (with $a=\sqrt k>0$, $b=-\sqrt k$) it has
exactly the two eigenvalues $\pm\sqrt k$. Conversely, by
Lemmas~\ref{lem:regular} and \ref{lem:dichotomy}(I), every connected
two-eigenvalue signed graph with $s=0$ has $A^2=kI$ with $\{0,\pm1\}$ entries
and zero diagonal, hence lies in $W(n,k)$. This is an exact correspondence.

The family is nonempty for infinitely many $(n,k)$. The unbalanced $4$-cycle
\[
A=\begin{pmatrix}0&1&0&-1\\ 1&0&1&0\\ 0&1&0&1\\ -1&0&1&0\end{pmatrix},
\qquad A^2=2I,
\]
is a connected member of $W(4,2)$ with eigenvalues $\pm\sqrt2$, each of
multiplicity $2$. Symmetric conference matrices give members of $W(n,n-1)$ on
the complete graph for $n\equiv2\pmod4$ (Paley type), and many further weights
$k$ occur; each connected such matrix yields a signed graph with spectrum
$\{\pm\sqrt k\}$.

\paragraph{Case (II), $s\ne0$.}
By Lemma~\ref{lem:dichotomy}(II), $a,b\in\mathbb Z$ and \eqref{eq:quad} holds
with integer $s\ne0$, $k=-ab$. Conversely, any connected $k$-regular
zero-diagonal $\{0,\pm1\}$ symmetric matrix $A$ with $A^2=sA+kI$,
$s\in\mathbb Z\setminus\{0\}$, $k=-ab$, satisfies $(A-aI)(A-bI)=0$ with $a\ne
b$ (as $k=-ab>0$), so by Lemma~\ref{lem:quad} it has exactly the two distinct
eigenvalues $a>b$. Thus case~(II) is precisely the class of connected integral
two-eigenvalue signed graphs.

\paragraph{Subcases (IIa),(IIb).}
By Proposition~\ref{prop:simple}, when one multiplicity equals $1$ the graph is
switching equivalent to the all-positive (resp.\ all-negative) $K_n$, with the
stated spectra. Conversely, the all-positive $K_n$ has $A=J-I$,
$A^2=(n-2)A+(n-1)I$, eigenvalues $n-1$ (simple) and $-1$; the all-negative
$K_n$ has $A=-(J-I)$, eigenvalues $1$ and $-(n-1)$ (simple). Both are connected
with exactly two distinct eigenvalues, realizing (IIa),(IIb) for all $n\ge2$.

\paragraph{Both multiplicities $\ge2$ in case (II).}
By Lemma~\ref{lem:regular} and Lemma~\ref{lem:dichotomy}(II), these are exactly
the connected $k$-regular zero-diagonal $\{0,\pm1\}$ matrices satisfying the
integer identity $A^2=sA+kI$ with $s\ne0$ in which neither extreme eigenvalue
is simple; equivalently, the connected integral two-eigenvalue signed graphs
not switching equivalent to a complete graph with a simple extreme eigenvalue.

\section{Conclusion}\label{sec:conclusion}

Combining Lemmas~\ref{lem:quad}, \ref{lem:regular}, \ref{lem:dichotomy} and
Proposition~\ref{prop:simple} with the converses of
Section~\ref{sec:examples} proves Theorem~\ref{thm:main}. A connected signed
graph has exactly two distinct eigenvalues if and only if it is $k$-regular and
its adjacency matrix satisfies \eqref{eq:quad}. These split, by the sign of
$s=a+b$, into the balanced family (case~I), exactly the connected symmetric
weighing matrices $W(n,k)$ with zero diagonal and spectrum $\{\pm\sqrt k\}$,
and the integral family (case~II), in which both eigenvalues are integers,
the multiplicity-one subcases being precisely the all-positive and
all-negative complete graphs up to switching, and the remaining members being
the connected $k$-regular signed graphs satisfying the integer identity
$A^2=sA+kI$ with $s\ne0$ and both multiplicities $\ge2$.
\qed

\end{document}
